\documentclass[lang=cn,11pt]{elegantpaper}
%放在最前面！！！

% ---------设置图片
\usepackage{graphicx}% 插入图片
\usepackage{subfigure}% subfigure——用来插入插入并列分布的图片% caption2——修改图注相关,这里好像不需要
\usepackage{float}% float——确定图片是否为浮动，而不是在一个固定的地方。
% \renewcommand{\figurename}{Fig.}
% \renewcommand{\captionlabeldelim}{ }%——这是修改分割符冒号为空格
% ---------设置代码块
\usepackage{listings}
% \usepackage{minted}%代码块高亮显示！！！

% ---------------------------------封面-------------------------------
\title{实验3-1  基于UDP设计可靠传输协议并编程实现}
\author{{1813800~~~~沈\ 哲}}
% \institute{{南开大学~~计算机学院,天津~~300000}}
\date{}% \version{0.07}
% % 如果想修改参考文献样式，请把这行注释掉% \usepackage[authoryear]{gbt7714}  % 国标
% % -------------------------------------------摘要和目录----------------------------------------
\begin{document}
    \maketitle
%     \begin{cnabstract}
%         沈哲，加油。
%         \par\textbf{关键词:~~}沈哲;加油
%     \end{cnabstract}
%     \begin{enabstract}
%         shen zhe
%         \par\textbf{Keywords:~~}shen;zhe
%     \end{enabstract}
    % \newpage
    \tableofcontents
    % \newpage
% -----------------------------------------------------------------------------------------
% ----------------------------------正文---------------------------------------------------
% -----------------------------------------------------------------------------------------
\section{实验要求}
        利用数据报套接字在用户空间实现面向连接的可靠数据传输，功能包括：建立连接、差错检测、确认重传。流量控制采用停等机制，完成给定测试文件的传输。

        要求实现单向传输。对于每一个任务要求给出详细的协议设计。完成给定测试文件的传输，显示传输时间和平均吞吐率。

\section{功能实现}
    \subsection{协议设计}
        UDP是传输层中面向无连接的协议，在编程上服务端和客户端是没有区别的，本实验实现从客户端（发送端）到服务端（接收端）的传输。

        基于UDP设计可靠传输协议，可以参考TCP协议设计相应的数据报的字段，本实验设计的数据报分为两部分——头部和数据部分。
        数据部分占65300字节。
        头部包括：序列号、确序号、检验和字段、标志位字段、长度字段。每一部分占2字节，共占10字节。
        
    \begin{lstlisting}
    Header:10 bytes
    | 2 bytes  | 2 bytes  | 2 bytes  | 2 bytes  |  2 bytes  |
    |   SEQ    |  ACKnum  | CheckSum |	  ID	|   Length  |

    ID:2 bytes
    |  FIN  |  ACK  |  SYN  |
    \end{lstlisting}
        
        SEQ是发送的序列号，主要在发送端使用，本次实验的停等协议中，只使用0和1两个序列号即可。

        ACKnum是确认序列号，主要在接收端使用，接收端回复的ACKnum等于发送端的SEQ。

        CheckSum是检验和字段，发送端制作数据包时计算检验和填入，接收端也对收到的数据包进行差错检验。

        ID是多位标志位，包括FIN（断开连接标志位）、ACK（确认标志位）、SYN（连接建立标志位）。


    \subsection{建立与断开连接}
    发送端发送只包含头部的包，SYN位置1，接收端接收到后发送ACK，连接建立。同理通过发送FIN与ACK进行“挥手”断开连接。

    \subsection{差错检测}
    采用UDP校验和计算方法。发送端在发送数据包时进行UDP校验和计算，结果写入CheckSum位。接收端接收到数据包后也进行校验和计算，结果为0xFFFF则无错误。
    \subsection{确认重传和停等协议}
    发送端收到确认序号ACKnum后再次发送下一个数据包。
    如果发送端发送的数据包丢失，接收端一直等待接收（程序里是在while循环中不断执行recvfrom函数），不发送数据包，或者接收端发送的数据包丢失，都会导致发送端收不到ACKnum，
    这时发送端就重传数据包。
     
    具体来说，发送端发送的数据包序列号只有0、1两种，如果发送的数据包SEQ=0,则期望收到ACKnum=0,之后再发送SEQ=1的数据包。若收不到ACKnum=0的数据包，就重传SEQ=0的数据包。

    对于接收端，收到某个序列号的数据包，发送确认序号，确认序号ACKnum等于序列号SEQ。如果收到重复的数据包（根据接收到包的序列号判断），就丢弃，但要发送ACKnum。

    这就实现了确认重传，能够处理发送端和接收端的丢包问题。

    \subsection{其他设计}
    采用非阻塞模式，设置一个超时时间（设置为10ms），超过此时间无响应就返回一个值（-1），这样编程更方便。经过测试，超时未响应recvfrom函数的返回值为-1，错误类型WSAGetLastError() 的返回值为 10060。
    \begin{lstlisting}
    // 设置发送超时
    setsockopt(socket_client,SOL_SOCKET,SO_SNDTIMEO, (char *)&nNetTimeout,sizeof(int));
    //设置接收超时
    setsockopt(socket_client,SOL_SOCKET,SO_RCVTIMEO, (char *)&nNetTimeout,sizeof(int));
    \end{lstlisting}


    利用随机数设置丢包，让发送端不发送数据包或者接收端不发送ACK数据包。测试传输的可靠性。



\section{代码分析(仅展示部分核心代码)}
    \subsection{公共头文件定义}
     \begin{lstlisting}
    //设置IP和端口号等
    #define server_Port 1001
    #define server_IP "127.0.0.1"
    #define client_Port 1002
    #define client_IP "127.0.0.1"
    //rand()%RAND_MOD_NUM 对随机数取模来决定发送还是丢弃数据包
    int RAND_MOD_NUM=15;
    //非阻塞模式的超时时间
    int nNetTimeout=10;//毫秒
    //缓冲区大小
    #define BUF_LEN 65310//比2**16小一点
    #define HEADER_LEN 10//头部长度
    #define DATA_LEN 65300
    /*  
    头部的设计
    Header:10 bytes
    | 2 bytes  | 2 bytes  | 2 bytes  | 2 bytes  |  2 bytes  |
    |   SEQ    |  ACKnum  | CheckSum |	  ID	|   Length  |
    
    ID:2 bytes
    |   FIN    |  ACK  |  SYN  |
    */
    struct HEADER
    {
    	short SEQ;//序列号
    	short ACKnum;//确认序号
    	short CheckSum;//检验和
    	short ID;//多个标志位
    	short Length;//数据长度，这是基本固定的
    };
    //各个标志位
    #define SYN 0x1//建立连接
    #define ACK 0x2//确认
    #define FIN 0x4//断开连接
    #define ID_bit 0x6
    #define ACKnum_bit 0x2
    #define SEQ_bit 0x0
    #define Length_bit 0x8
\end{lstlisting}

    \subsection{发送端和接收端都使用的函数}
    \subsubsection{makePackage}
    用于制作数据包的头部，发送端需要设置校验和、发送序列号、数据长度。接收端需要设置确认序号。
    \begin{lstlisting}
    header.Length = length;
    header.SEQ =(1+header.SEQ)%2;//0 1 0 1 
    header.CheckSum = 0;
    memcpy(&sendbuf, &header, HEADER_LEN);
    \end{lstlisting}
    注意建立连接时填充SYN字段，结束连接时填充FIN字段。发送端在读取数据后，需要再次填充CheckSum字段。
    \subsubsection{doCheckSum}
    计算UDP校验和的函数。发送端计算校验和填入CheckSum字段。接收端计算校验和判断是否等于0xFFFF，不等于0xFFFF则数据出错。
    \begin{lstlisting}
    //先将char型发送缓冲区转换成为16位的unsigned short型数据
    //需要分情况，考虑缓冲区长度
    array = new unsigned short[count];
    for (int i = 0; i < count; i++)
    {
        memcpy(&array[i], &sendbuf[i * 2], 2);
    }
    //计算校验和
    while (count--)
    {
        sum += array[i++];
        sum = (sum >> 16) + (sum & 0xFFFF);
    }
    header.CheckSum=~sum;
    memcpy(&sendbuf[0],&header, 10);
    \end{lstlisting}
    \subsubsection{CheckSEQ}
    发送端检查接收的ACKnum是否等于发送序列号。接收端检查接受的是否是新的序列号的数据包（若不是需要丢弃）。
    \begin{lstlisting}
    memcpy(&seq, &recvbuf[ACKnum_bit], 2);//ACKnum字段
    if (seq == header.SEQ)
        return true;
    \end{lstlisting}
    
    \subsubsection{IsACK,IsFIN,IsSYN}
    检查这些标志位是否置位。

    \subsection{客户端（发送端）}
    先是WinSocket的初始化，创建socket，设置超时时间等操作。
    \begin{lstlisting}
    // 设置发送超时
    setsockopt(socket_client,SOL_SOCKET,SO_SNDTIMEO, (char *)&nNetTimeout,sizeof(int));
    //设置接收超时
    setsockopt(socket_client,SOL_SOCKET,SO_RCVTIMEO, (char *)&nNetTimeout,sizeof(int));
    \end{lstlisting}

    发送包含SYN的包，准备建立连接。

    连接建立后，先发送文件名。之后开始读文件（每次读取数据缓冲区大小的数据量）。制作数据包（包括头部的各个字段，SEQ进行加1模2操作实现01交替，检待装入数据后计算校验和）。
    发送数据包，这里随机丢弃进行测试。
    \begin{lstlisting}
    file.read(&sendbuf[10], DATA_LEN);
    makePackage(file.gcount());
    doCheckSum(file.gcount());
    //发送数据包,随机丢弃（不发送）
    if(rand()%RAND_MOD_NUM!=0)//遇到rand()%15==0丢弃
        sendto(socket_client, sendbuf, BUF_LEN, 0, (SOCKADDR*)&server_addr, sizeof(SOCKADDR));
    else
        cout<<"丢弃数据包！！！！！！！！！！！"<<endl;
    \end{lstlisting}

    接收到ACK，并且确认序列号正确。
    \begin{lstlisting}
    cout << "第" << trans_num << "次传输成功！" <<endl<< "传输" << file.gcount()<<"字节" << endl;
    trans_bytes+=file.gcount();
    memset(&sendbuf[10], 0, DATA_LEN);//清空sendbuf
    continue;
    \end{lstlisting}

    未接收到ACK，进行超时重传。
    \begin{lstlisting}
    //接收ACK
    recv_Ret = recvfrom(socket_client, recvbuf, HEADER_LEN, 0, (SOCKADDR*)&server_addr, &recv_para_len);
    if (recv_Ret < 0)
    { //(WSAGetLastError() == 10060)
        cout << "超时，未接收到ACK，重传数据包！！！" << endl;
        /*处理：
        1.读文件指针回退，
        2.trans_num--
        3.接收端做丢弃acknum，就要处理来自发送端的重复的包，就必须检查序列号，发送端就要做好0 1交替
        调用两次makePackage即可！！！
        刚刚发了包0，多调用一次makePackage（SEQ==1）,再发包0时就可以了makePackage（SEQ==0）
        */
        file.seekg(-file.gcount(),ios::cur);//向前移动
        trans_num--;
        makePackage(0);//重传的处理，多调用一次使得序列号和之前一样
        continue;
    }
    \end{lstlisting}
    \subsection{服务端（接收端）}
    开始处理和发送端相同。建立连接后接收文件名。之后开始接收数据包，写入文件，发送ACK确认数据包（也进行随机丢弃）。
   
    CheckSum函数进行校验和计算，CheckSEQ确保接收到的是下一个序列的数据包，否则接收重复丢弃即可，IsFIN判断是否传输完成。
   
    接收数据包并写入文件。
   \begin{lstlisting}
    file.write(&recvbuf[10], DATA_LEN); 
    memcpy(&data_len, &recvbuf[Length_bit], 2);
    trans_bytes += data_len;
    memset(&recvbuf[10], 0, DATA_LEN); //及时清空缓存区
    makePackage();
    if(rand()%RAND_MOD_NUM!=0)
    {
        sendto(socket_server, sendbuf, HEADER_LEN, 0, (SOCKADDR *)&client_addr, sizeof(SOCKADDR));
    }
    else
    {
        cout<<"丢弃ACK包！！！！！！！！！！！"<<endl;
    }
    continue;
    \end{lstlisting}

    接收到重复的数据包，不写入文件，丢弃即可。
    \begin{lstlisting}
    //本该收包1（然后回复acknum==1），结果收到包0，
    //由于刚刚收到的包0，此时header.acknum==0，正常调用makePackage会使得header.acknum==1
    //但是这时候还要回复acknum==0，函数细节-->memcpy(&header.ACKnum, &recvbuf[SEQ_bit], 2);
    //还是直接用makePackage即可，因为是根据收到的recvbuf[SEQ_bit]确定acknum，
    //不像发送端的SEQ是每次调用makePackage会01交替
    cout<<"接收到重传的数据包，丢弃，发送ACK！！！"<<endl;
    makePackage();
    sendto(socket_server, sendbuf, HEADER_LEN, 0, (SOCKADDR *)&client_addr, sizeof(SOCKADDR));
    memset(&recvbuf[10], 0, DATA_LEN); //及时清空缓存区
    continue;
    \end{lstlisting}

\section{结果展示}
程序可以连续发送任意类型的文件，之后也可以设置用户输入IP地址、端口号，文件名，选择设置丢包率（随机数设置）、超时时间等等，能够实现文件的无损正确传输。
超时时间越短（这里是10ms），缓冲区越大（这里是比2的16次方小一些，不能超过，取65310），文件发送速度越快。
\subsection{发送端丢弃数据包}
\begin{figure}[H] %figure环境，h默认参数是可以浮动，不是固定在当前位置。如果要不浮动，你就可以使用大写float宏包的H参数，固定图片在当前位置，禁止浮动。
    \centering %居中
    \includegraphics[width=0.8\textwidth]{image/1.png}  %1:设置图片充满文档的大小
    \label{png1} %添加标签，从1递增
\end{figure}

\subsection{接收端丢弃ACK包}
\begin{figure}[H] %figure环境，h默认参数是可以浮动，不是固定在当前位置。如果要不浮动，你就可以使用大写float宏包的H参数，固定图片在当前位置，禁止浮动。
    \centering %居中
    \includegraphics[width=0.8\textwidth]{image/2.png}  %1:设置图片充满文档的大小
    \label{png2} %添加标签，从1递增
\end{figure}

\subsection{整体效果}
\begin{figure}[H] %figure环境，h默认参数是可以浮动，不是固定在当前位置。如果要不浮动，你就可以使用大写float宏包的H参数，固定图片在当前位置，禁止浮动。
    \centering %居中
    \includegraphics[width=0.8\textwidth]{image/3.png}  %1:设置图片充满文档的大小
    \label{png2} %添加标签，从1递增
\end{figure}
\section{总结反思}
\subsection{巩固了C++基础知识}
在实验过程中因代码不严谨出现许多错误，通过调试更加熟悉了字符串函数的使用（如memcpy），文件读写操作等等。
\subsection{更熟悉了UDP和TCP协议以及基于winsock的编程}
UDP是传输层中面向无连接的协议，本实验以UDP为平台，进行可靠传输协议的设计，还是能够参考TCP的一些思想，对两种协议的了解也更深入。
在设计思考过程中，复习了rdt不同版本的协议，对可靠传输协议的设计更加熟悉。




% -------------------------------参考文献------------------------------------
\end{document}




            
  
